\section{Diccionario razonado de las variables que usa el análisis}
\label{sec:anexo-diccionario}
La Tabla~\ref{tab:diccionario} lista las variables que entran al análisis. El diccionario completo del MC4 (76
variables) está en el repositorio (\texttt{outputs/tablas/diccionario\_razonado.csv}). \emph{Estado}: observado
(verificado en los datos), inferido (interpretación nuestra) o supuesto (sin confirmación posible; Anexo~\ref{sec:anexo-supuestos}).

\begin{table}[!h]
\centering
\caption{Diccionario razonado. Fuente: MC4 (archivos MC4\_PSQ) o SKP (archivos SKP\_PSQ). Tipo: especificación (se fija
antes de producir), medición, etiqueta (la calcula el sistema o la asigna un inspector) o identificador.}
\label{tab:diccionario}
\footnotesize
\begin{tabular}{L{3.2cm}cL{4.4cm}cL{1.8cm}L{2.5cm}c}
\toprule
Variable (columna) & Fuente & Qué es & Unidad & Tipo & Papel & Estado \\
\midrule
Material Salida / Material Entrada & MC4 / SKP & Identificador del rollo; llave de la unión & --- & identificador & llave & observado \\
Programa & SKP & Tanda de rollos procesados de corrido en el Skinpass (mediana \MedianaRollosPrograma) & --- & identificador & grupo, agrupación de errores & inferido \\
Rodillo de Trabajo & SKP & Rodillo de trabajo montado (\NRodillos{} rodillos) & --- & identificador & alternativa al programa (R9) & observado \\
Norma, Grado & MC4 & Estándar del acero y su variante; juntos son el producto & --- & especificación & confusor & observado \\
Espesor y Ancho Objetivo & MC4 & Dimensión pedida, en rangos & mm & especificación & confusor & observado \\
Fecha Inicio & MC4 & Inicio de laminación; da el mes y el corte de censura & --- & medición & confusor (periodo) & observado \\
Corona Salida, Espesor & MC4 & Corona de salida y espesor medido; su cociente es la corona relativa & µm, mm & medición & exposición & supuesto (S6) \\
Temp Enrollado Real y Obj & MC4 & Temperatura de enrollado real y objetivo & °C & medición / esp. & exposición secundaria & observado \\
Defecto (i) = INCLUSIONES & MC4 & Defecto con origen en el acero & --- & etiqueta & control negativo & observado \\
Acuñamiento Salida & MC4 & Cuña de salida & µm & medición & comprobación (R7) & supuesto \\
Preset Bending & SKP & Bending fijado por la receta & kN & especificación & componente de $\Y$ & supuesto (S1, S2) \\
Bending Prom. Real [kN] & SKP & Bending real promedio aplicado por el lazo & kN & medición & componente de $\Y$ & observado \\
Fuerza, Tensiones reales & SKP & Fuerza de laminación y tensiones promedio & kN & medición & simultáneas (no entran) & observado \\
Elongación, Modo Molino & SKP & Elongación lograda; modo fuerza o elongación & \%, --- & medición & indicador secundario & observado \\
Registro Automatico (i), Linea Origen (i) & SKP & Si el defecto lo puso el sistema y en qué línea nació & --- & etiqueta & define el indicador de inspector & observado \\
Dictamen & SKP & ACEPTADO, RETENIDO o REPROCESO & --- & etiqueta & indicador secundario & observado \\
Peso & SKP & Peso de cada pieza; pondera la agregación por rollo & kg & medición & ponderador & observado \\
\bottomrule
\end{tabular}
\end{table}

\section{Preguntas al Socio Formador y supuestos}
\label{sec:anexo-supuestos}
En la Entrega~I, Ternium respondió dos preguntas: el acuñamiento de un solo lado no es una prioridad operativa (por
eso salió de la pregunta) y compartió señales de bending del Skinpass de 2026 a 25~Hz, que no coinciden con ningún
rollo de 2025 y quedaron fuera de alcance. Las demás preguntas (qué resume cada valor, unidades, registro de rodillos,
reglas automáticas) ya no pueden consultarse; cada una se convirtió en un supuesto explícito con su evidencia y su
prueba de sensibilidad (Tabla~\ref{tab:supuestos}).

\begin{table}[!h]
\centering
\caption{Supuestos, evidencia en los datos y prueba de sensibilidad.}
\label{tab:supuestos}
\footnotesize
\begin{tabular}{cL{3.5cm}L{6.6cm}L{2.9cm}L{2.6cm}}
\toprule
& Supuesto & Evidencia & Si es falso & Sensibilidad \\
\midrule
S1 & El preset está en kN, como el bending real & Misma escala (percentiles 1--99: preset \PresetPuno{} a \PresetPnoventaynueve; real \RealPuno{} a \RealPnoventaynueve~kN); el real sigue al preset (pendiente \PendRealPreset, $r=\CorrRealPreset$) & $\Y$ mezclaría unidades & R8; el bending real (sin preset) da la misma pendiente \\
S2 & Preset cero significa ``sin preset'' & \PctPresetCero\,\% de los rollos; con preset cero el bending real sigue activo (mediana \BendRealPresetCero~kN; exactamente cero en \PctRealCeroPresetCero\,\%) & $\Y$ sería otra cosa en esos rollos & R5 \\
S3 & El preset se fija antes del rollo y no reacciona a él & Único en \PctPresetUnicoTodos\,\% de los grupos programa $\times$ producto $\times$ dimensión (\PctPresetUnicoDos\,\% con dos o más rollos); no responde a la corona (\PendPreset~kN/pp) & El preset sería una respuesta & Se describe; afecta la lectura \\
S4 & El programa agrupa rollos con el mismo rodillo de trabajo & \PctProgUnRodillo\,\% de los programas usa un solo rodillo; mediana de \MedianaProgRodillo{} programas por rodillo & El programa no reflejaría el rodillo & R9 \\
S5 & El bending real responde a un lazo de planeza & El real se aparta del preset (mediana de $\Y$: \YMediana~kN; $|\Y|>50$~kN en \PctYMayorCincuenta\,\%) & El bending sería una decisión manual & Se describe en el diagrama \\
S6 & La corona está en µm & Mediana \CoronaMedianaUm~µm (p5--p95: \CoronaPcincoUm--\CoronaPnoventaycincoUm) con espesor mediano de \EspesorMediano~mm: corona relativa de \CoronaRelMediana\,\% & Cambia la escala de la pendiente, no su signo & Pendiente también por DE \\
\bottomrule
\end{tabular}
\end{table}

\section{Tabla completa de contrastes examinados}
\label{sec:anexo-contrastes}
Examinamos \NContrastes{} contrastes. Cada uno, con su estimación, intervalo, unidades, valor $p$, valor $q$ de
Benjamini--Hochberg cuando corresponde, rollos y grupos, está en \texttt{outputs/tablas/contrastes\_completos.csv}. La
Tabla~\ref{tab:contrastes} los resume por familia; en el cuerpo del artículo solo aparecen los que contestan la pregunta.

\begin{table}[!h]
\centering
\caption{Resumen de los contrastes por familia. Mín.\ y Máx.: rango de las estimaciones en la familia. IC excluye 0:
cuántos intervalos de 95\,\% excluyen el cero. $q_{\min}$: menor valor $q$ de Benjamini--Hochberg.}
\label{tab:contrastes}
\footnotesize
\begin{tabular}{L{4.6cm}rrrlrL{4.2cm}}
\toprule
Familia & $n$ & Mín. & Máx. & IC excluye 0 & $q_{\min}$ & Unidades \\
\midrule
8.B descomposición (orden del plan) & 4 & 0.22 & 22.07 & 3 de 4 & -- & \% de la SC de Y \\
8.B descomposición (programa primero) & 4 & 0.22 & 26.29 & 3 de 4 & -- & \% de la SC de Y \\
V4 programa vs. periodo & 4 & 13.23 & 22.07 & 4 de 4 & -- & \% de la SC de Y \\
8.C asociación principal & 2 & −3.90 & −0.23 & 1 de 2 & -- & kN por unidad \\
Indicadores secundarios & 10 & −3.78 & 0.12 & 2 de 10 & -- & por unidad de la exposición \\
8.D heterogeneidad por producto & 35 & −30.04 & 10.24 & 1 de 35 & 0.70 & kN por pp \\
V1 control negativo & 2 & −0.27 & 0.09 & 0 de 2 & -- & kN \\
V3 robustez: pendiente de corona & 16 & −5.71 & −2.18 & 15 de 16 & -- & kN por pp \\
V3 robustez: exceso del grupo − producto & 16 & 4.34 & 24.73 & -- & -- & puntos de \% de la SC \\
10 confusión no medida & 1 & 1.47 & 1.47 & -- & -- & \% de varianza residual \\
8.E predictivo & 6 & −2.21 & 0.84 & -- & -- & kN (media entre particiones) \\
12 experimento: tamaño del efecto & 4 & 0.82 & 8.04 & -- & -- & kN (media de 5 particiones agrupadas) \\
\bottomrule
\end{tabular}

\end{table}

\section{Plan de análisis}
\label{sec:anexo-plan}
El plan (\texttt{docs/plan\_de\_analisis.md}) fijó, antes de la estimación principal, la pregunta y la afirmación más fuerte
que esperábamos sostener, la unión y la población con sus exclusiones, el resultado y los indicadores secundarios, el
conjunto de ajuste derivado del diagrama, el diagnóstico de soporte, los análisis, los cuatro anclajes de validación
con sus criterios de refutación, las nueve alternativas de robustez y los seis supuestos. Su bitácora registra, con fecha
y motivo, los cambios posteriores. Los principales: el intervalo del exceso sobre el azar se calcula con mitades de
programas en lugar de bootstrap con reemplazo (el bootstrap se conserva en el repositorio); se definió de antemano qué
cuenta como mejora clara del modelo predictivo; el signo físico esperado se presenta como argumento mecánico del equipo,
porque la referencia prevista no lo confirma; y la propuesta de experimento, escrita después de ver los resultados,
agrega el factor de recentrado por producto y cambia la unidad de asignación del programa al rollo dentro de programa.
